\documentclass[10pt,a4paper]{amsart}
\usepackage[margin=19mm]{geometry}
\usepackage{amsmath,amssymb,mathtools}
\usepackage[scr=rsfs,cal=euler]{mathalfa}
\usepackage{xfrac}
\usepackage[unicode,hidelinks,bookmarks=false]{hyperref}
\newcommand{\ie}{i.e.\ }
\newcommand{\cf}{cf.\ }
\newcommand{\resp}{respectively\ }
\newcommand{\Subsection}{\subsection}
\providecommand{\nobreakdash}{\nobreak\hbox{-}}
\setlength{\emergencystretch}{2em}
\multlinegap0pt
\usepackage{bm}
\usepackage{bbm}
\usepackage{dsfont}
\usepackage{xspace}
\usepackage{cancel}
\usepackage{mathtools}
\usepackage{amsthm}
\makeatletter
\@ifundefined{theorem}{\newtheorem{theorem}{Theorem}}{}
\@ifundefined{lemma}{\newtheorem{lemma}[theorem]{Lemma}}{}
\makeatother
\newcommand{\df}[1]{\ensuremath{\mathrm{df}({#1})}\xspace}
\title[Cubic graphs of colouring defect 3 and conjectures of Berge ]{Cubic graphs of colouring defect 3 and conjectures of Berge and Alon-Tarsi}
\author{J\'an Karab\'a\v{s}, Edita M\'a\v{c}ajov\'a, Roman Nedela, Martin \v{S}koviera}
\date{Source: arXiv:2505.17569v1 (CC BY 4.0)}
\begin{document}
\maketitle

\noindent\emph{Source and licence.} Cubic graphs of colouring defect 3 and conjectures of Berge and Alon-Tarsi. Version arXiv:2505.17569v1 (CC BY 4.0), distributed under the Creative Commons Attribution 4.0 International licence, as recorded in the arXiv licence metadata for this exact version and shown on the abstract page https://arxiv.org/abs/2505.17569v1. Full rights notice: licenses/arxiv-2505.17569-CC-BY-4.0.txt.\par\smallskip

\noindent\textit{Editorial scope.} Counted unit: the Theorem whose source label is \texttt{thm:main} in the article (source0.tex lines 1777--1796, source label \texttt{thm:main}), delivered together with the article's own complete original proof of it (source0.tex lines 1798--1960, one \texttt{proof} environment of 163 lines). No mathematics is written, repaired or re-derived: statement and proof are the article's own text line for line. Required context retained uncounted: thm:decomp-def3 (lines 1106--1124); lem:2+3-sum-4cover (lines 1228--1233); thm:triancore (lines 1306--1309); thm:pmi-c4c (lines 1507--1512); lem:2sum (lines 1563--1567); lem:pmi4sum3gen (lines 1609--1615); thm:3-sumpmi5 (lines 1668--1683); thm:3-sumpmi5-converse (lines 1704--1710). Every definition, display and label of those lines is retained unchanged. Omitted from this package and not redistributed: 1--1105, 1125--1227, 1234--1305, 1310--1506, 1513--1562, 1568--1608, 1616--1667, 1684--1703, 1711--1776, 1797--1797, 1961--2382. Nothing inside a retained excerpt is omitted or altered: every retained body is delivered exactly as the article's lines are written, so no omission receipt of any kind accompanies this package. Citations: the retained excerpts carry no citation command at all, so no bibliography entry of the article is transcribed and no reference is redistributed. Reference adaptation: none was needed and none was made; every reference inside the delivered unit resolves inside it, so no unresolved reference or citation is produced. Printed slips kept literal: no printed word, symbol, hypothesis or number of the counted statement or of its proof was altered. Layout and class: the article's own class is \texttt{amsart} with packages including color, graphicx, subfigure, float, url, enumerate, xspace, multirow, hhline, lmodern, caption; the wrapper is the lane's pinned \texttt{amsart} editorial wrapper at 10pt on A4, which restates the theorem-like environments the retained text needs and adds the article's own macro definitions that the retained text needs (\texttt{\textbackslash df}), with extra wrapper packages (bm, bbm, dsfont, xspace, cancel, mathtools). The delivered numbering is the unit's own. Assets: no retained span contains a figure environment, a graphics-inclusion command, a PDF-page inclusion, a table or a TikZ picture; no image file is copied into this package, the package declares no asset dependency and no image file is redistributed with it. Licence: version arXiv:2505.17569v1 of this article is distributed under the Creative Commons Attribution 4.0 International licence, as recorded in the arXiv licence metadata for this exact version and shown on the abstract page https://arxiv.org/abs/2505.17569v1. A full rights notice is delivered at licenses/arxiv-2505.17569-CC-BY-4.0.txt. Characters: every retained excerpt is pure ASCII, so the delivered engine, XeLaTeX with its default Latin Modern text font, reports no missing character.
\begin{theorem}\label{thm:decomp-def3}
Let $G$ be a $2$-connected cubic graph of defect $3$ and let
$C$ be a hexagonal core of $G$. Then $G$ admits a decomposition
$\mathcal{D}=\{G_1,\dots,G_m\}$ such that all decomposition
factors are cyclically $4$-edge-connected, except possibly
$G_1$, and the following hold:
\begin{itemize}
\item[{\rm (i)}] $G_1$ has defect $3$ and $C$ is a
    hexagonal core of $G_1$;
\item[{\rm (ii)}] each $G_i$ with $i\ge 2$ is
    $3$-edge-colourable; and
\item[{\rm (iii)}] if $G_1$ is not cyclically
    $4$-edge-connected, then it contains a triangle $T$
    such that $C$ intersects $T$ and $G_1/T$ is cyclically
    $4$-edge-connected.
\end{itemize}
Moreover, $G$ can be reconstructed from $\mathcal{D}$ by
repeated application of $2$-sums and $3$-sums.
\end{theorem}

\begin{lemma}\label{lem:2+3-sum-4cover}
Let $H$ and $K$ be $2$-connected cubic graphs. If $\pi(H)=4$
and $K$ is $3$-edge-colourable, then $\pi(H\oplus_2 K)=4$ and
$\pi(H\oplus_3 K)=4$ for any choice of distinguished edges or
vertices.
\end{lemma}

\begin{theorem}\label{thm:triancore}
Let $G$ be a snark with colouring defect $3$. If $G$ contains a hexagonal
core intersecting a triangle, then $\pi(G)=4$.
\end{theorem}

\begin{theorem}\label{thm:pmi-c4c}
Every snark $G$ of defect $3$ satisfies Berge's conjecture,
that is, $\pi(G)\le 5$. Moreover, if $G$ is cyclically
$4$-edge-connected, then $\pi(G)=4$ unless $G$ is the Petersen
graph, which has $\pi(Pg)=5$.
\end{theorem}

\begin{lemma}\label{lem:2sum}
Let $G=H\oplus_2 K$ be a $2$-sum of $2$-connected cubic graphs
$H$ and $K$, where $K$ is $3$-edge-colourable. Then $\pi(G)\ge
5$ if and only if $\pi(H)\ge 5$.
\end{lemma}

\begin{lemma}\label{lem:pmi4sum3gen}
Let $G=H\oplus K$ be a $2$-sum or a $3$-sum of $2$-connected
cubic graphs $H$ and $K$ where $H$ is a snark with $\pi(H)\ge
5$ and $K$ is $3$-edge-colourable. If $u$ is an apex of $H$
different from the distinguished vertex for the $3$-sum, then
$u$ is an apex of $G$.
\end{lemma}

\begin{theorem}\label{thm:3-sumpmi5}
Let $H$ and $K$ be $2$-connected cubic graphs, where $H$ is a
snark with $\pi(H)\ge 5$ and $K$ is $3$-edge-colourable
quasi-bipartite graph. If $G=H\oplus_3 K$ is a $3$-sum with
distinguished vertices $u$ and $v$, where
%$u$ is an apex of $H$ and
$v$ forms a trivial component of the quasi-partite set of $K$,
then the following hold:
\begin{itemize}
\item[{\rm (i)}] $\pi(G)\ge 5$.

\item[{\rm (ii)}] $G$ is quasi-bipartite with bipartising
    set inherited from $K$, and $H-u$ is an element of the
    quasi-partite set of $G$ replacing $v$.
\end{itemize}
\end{theorem}

\begin{theorem}\label{thm:3-sumpmi5-converse}
Let $G=H\oplus_3 K$ be a proper $3$-sum of $2$-connected cubic
graphs where $\pi(H)\ge 5$, $K$ is $3$-edge-colourable, and the
distinguished vertex of $H$ is an apex. If $\pi(G)\ge 5$, then
$K$ is quasi-bipartite and its distinguished vertex forms a
trivial component of the quasi-partite set of $K$.
\end{theorem}

\begin{theorem}\label{thm:main}
Every $2$-connected cubic graph with colouring defect $3$ and
perfect matching index $5$ is either isomorphic to the Petersen
graph $Pg$ or arises from the Petersen graph with a fixed
$6$-cycle $C$ by combining the following two operations:
\begin{itemize}
\item[{\rm (O1)}] insertion of connected
    $3$-edge-colourable edge-deleted cubic graphs into any
    number of edges from $E(Pg)-E(C)$; and
\item[{\rm (O2)}] correct substitution of any number of
    vertices from $V(Pg)-V(C)$ with connected
    vertex-deleted $3$-edge-colourable quasi-bipartite
    cubic graphs.
\end{itemize}
Conversely, every cubic graph constructed from the Petersen
graph by means of operations {\rm (O1)} and {\rm (O2)} has
defect $3$ and perfect matching index~$5$. Moreover, the
$6$-cycle $C$ constitutes a hexagonal core of the resulting
graph.
\end{theorem}

\begin{proof}
$(\Leftarrow)$ We show that every graph $G$ constructed from
$Pg$ with a fixed $6$-cycle $C$ by means of operations (O1) and
(O2) has defect $3$ and perfect matching index $5$. As regards
its perfect matching index, Theorem~\ref{thm:pmi-c4c} implies
that it is enough to show that $\pi(G)\ge 5$. Thus it
sufficient to prove the following by induction on $t$:

\medskip\noindent
Claim 1. \emph{Let $G_t$ be the graph obtained from the
Petersen graph with a fixed $6$-cycle $C$ by applying $t$
operations {\rm (O1)} and {\rm (O2)} in any order. Then
$\pi(G_t)\ge 5$ and $G_t$ has a $3$-array $\mathcal{M}_t$ whose
core is $C$.}

\medskip\noindent
\emph{Proof of Claim 1.} Before starting the induction
procedure it may be useful to realise that $0\le t\le 13$. For
$t=0$ the statement is obviously true. Assume that the
statement is true for some $t\ge 0$, and consider the graph
$G_{t+1}$.

First assume that $G_{t+1}$ arises from $G_t$ by applying an
edge-insertion into an edge $e\in E(Pg)-E(C)$. Thus there
exists a connected $3$-edge-colourable graph $K$ and an edge
$k$ of $K$ such that $G_{t+1}=G_t\oplus_2 K$ with $e$ and $k$
being the distinguished edges for the $2$-sum, respectively.
Let $R$ be the principal $2$-edge-cut of the $2$-sum. Due to
the its choice, $e$ is simply covered with respect to the
$3$-array $\mathcal{M}_t=\{M_1^t,M_2^t,M_3^t\}$. Without loss
of generality we may assume that $e\in M_1^t$. Let $N_1$, $N_2$
and $N_3$ be the colour classes of an arbitrary
$3$-edge-colouring of~$K$; we can clearly assume that $k\in
N_1$. Set $M_1^{t+1}=(M_1^t-\{e\})\cup (N_1-\{k\})\cup R$ and
$M_i^{t+1}=M_i^t\cup N_i$ for $i\in\{2,3\}$. Obviously,
$\mathcal{M}_{t+1}=\{M_1^{t+1},M_2^{t+1},M_3^{t+1}\}$ a
$3$-array for $G_{t+1}$ with $C$ being its hexagonal core. The
fact that $\pi(G_{t+1})\ge 5$ follows immediately from
Lemma~\ref{lem:2sum}.

Next assume that $G_{t+1}$ arises from $G_t$ by substituting a
vertex $u\in V(Pg)-V(C)$. Now, there is a connected
$3$-edge-colourable quasi-bipartite cubic graph $K$ and a
vertex $v$ of $K$ such that $G_{t+1}=G_t\oplus_3 K$, with $u$
and $v$ being the distinguished vertices for the $3$-sum.
Theorem~\ref{thm:3-sumpmi5} implies that $\pi(G)\ge 5$.
Moreover, since $u$ does not lie on $C$, each edge incident
with $u$ is simply covered with respect to the $3$-array
$\mathcal{M}_t$. Since $K$ is $3$-edge-colourable, each perfect
matching $M_i^t\in\mathcal{M}_t$ can be extended to a perfect
matching $M_i^{t+1}$ of $G_{t+1}$ in such a way that the core
of the $3$-array
$\mathcal{M}_{t+1}=\{M_1^{t+1},M_2^{t+1},M_3^{t+1}\}$ is
again~$C$.

$(\Rightarrow)$ Let $G$ be a $2$-connected cubic graph with
$\df{G}=3$ and $\pi(G)=5$ different from the Petersen graph
$Pg$. Let $C$ be an arbitrary hexagonal core of $G$. We intend
to show that $G$ can be obtained from $Pg$ by a series of
operations (O1) and (O2) leaving $C$ intact.  We start by
applying Theorem~\ref{thm:decomp-def3}. It states that $G$
admits a decomposition $\mathcal{G}=\{G_1,\ldots,G_m\}$ such
that $G_1$ has defect $3$ and inherits $C$ as a hexagonal core,
all other decomposition factors are $3$-edge-colourable and
cyclically $4$-edge-connected, while $G_1$ is either cyclically
$4$-edge-connected or contains a triangle $T$ intersected by
$C$ such that $G_1/T$ is cyclically $4$-edge-connected.
Moreover, $G$ can be reconstructed from $\mathcal{G}$ by a
repeated application of $2$-sums and proper $3$-sums. We first
prove that that $\pi(G_1)\ge 5$. Suppose to the contrary that
$\pi(G)=4$. Since $G$ can be obtained from $\mathcal{G}$ by
applying $2$-sums and $3$-sums, we can repeatedly apply
Lemma~\ref{lem:2+3-sum-4cover} to deduce that $\pi(G)=4$, which
contradicts the assumption. Thus $\pi(G_1)\ge 5$. 
We further claim that $G_1$ is
cyclically $4$-edge-connected, implying that $\mathcal{G}$ is
the canonical decomposition of $G$. If $G_1$ is not cyclically
$4$-edge-connected, then, by Theorem~\ref{thm:decomp-def3}, it
contains a triangle $T$ intersecting~$C$, which is s a
hexagonal core of $G_1$. Theorem~\ref{thm:triancore} now yields
that $\pi(G_1)=4$, which again is a contradiction. Therefore
$G_1$ is a cyclically $4$-edge-connected cubic graph with
$\df{G_1}=3$ and $\pi(G_1)\ge 5$; according to
Theorem~\ref{thm:pmi-c4c}, $G_1$ is the Petersen graph $Pg$.

To show that $G$ can be obtained from the Petersen graph by
using operations (O1) and (O2) we build $G$ from $G_1=Pg$ by
using the canonical decomposition $\mathcal{G}$ of $G$. Recall
that $C$ is a hexagonal core of $G_1$. Further recall that
certain edges and vertices in $G_1$ are distinguished for
$2$-sums and $3$-sums applied during the reconstruction of $G$
from~$\mathcal{G}$. The same holds for the remaining
decomposition factors as well. The uniqueness of the
decomposition implies that during the reconstruction process
$2$-sums and proper $3$-sums can be used in any order. Thus we
can begin with $2$-sums and $3$-sums that do not use
distinguished edges or vertices from~$G_1$. Applying $2$-sums
and $3$-sums to all such distinguished edges and vertices
results in a decomposition $\mathcal{Q}=\{Q_1,\ldots, Q_r\}$
of~$G$, where $Q_1=G_1$ is the Petersen graph and each $Q_i$
with $i\ge 2$ is a connected $3$-edge-colourable graph. Now we
reconstruct $G$ from $\mathcal{Q}$. It is clear that each
$2$-sum or $3$-sum applied in this reconstruction involves an
edge or a vertex of $Q_1=Pg$. In fact, it follows from the
definition of the decomposition of a cubic graph along a small
independent edge cut that each $Q_i$ with $i\ge 2$ contains
exactly one distinguished element whose counterpart lies
in~$Q_1$. The distinguished vertices and edges lying in $Q_1$
must occur outside $C$, because $C$ has been left intact during
the entire decomposition process. In particular, $r\le 14$.

We reorder $\mathcal{Q}$ in such a way that the graphs $Q_2,
\ldots, Q_p$ have distinguished edges and the graphs $Q_{p+1},
\ldots, Q_r$ have distinguished vertices. We can now construct
$G$ inductively by setting $H_1=Q_1$ and $H_{i+1}=H_i\oplus_2
Q_{i+1}$ for $i\in\{1,\ldots, p-1\}$ and $H_{i+1}=H_i\oplus_3
Q_{i+1}$ for $i\in\{p, \ldots, r-1\}$. Consequently, $G=H_r$.
Note that each $H_i$, with $p\le i\le r-1$, has a distinguished
vertex, denoted by $u_i$, which is an original vertex of
$Q_1=Pg$.

What remains to be proved is the following statement.

\medskip\noindent
Claim 2. \emph{For every $i\ge p+1$ the graph
$Q_i\in\mathcal{Q}$ is quasi-bipartite and each $H_i$ arises
from $H_{i-1}$ by a correct substitution of its distinguished
vertex $u_i$.}

\medskip\noindent
\emph{Proof of Claim 2.} Our aim is to apply
Theorem~\ref{thm:3-sumpmi5-converse}. To this end we first
prove that every graph $H_i$ with $i\ge p$ has $\pi(H_i)\ge 5$
and $u_i$ is an apex of $H_i$. Recall that each $Q_i$ with
$2\le i\le p$ is $3$-edge-colourable, so we can repeatedly
apply Lemma~\ref{lem:2sum} to conclude that $\pi(H_i)\ge 5$
whenever $1\le i \le p$. In particular, $\pi(H_p)\ge 5$. To
prove that $\pi(H_i)\ge 5$ also for each $i\ge p+1$ we start
with the fact that $G=H_r=H_{r-1}\oplus_3 Q_r$ where
$\pi(H_r)=\pi(G)=5$ and $Q_r$ is $3$-edge-colourable. It
follows that $\pi(H_{r-1})\ge 5$, because otherwise
Lemma~\ref{lem:2+3-sum-4cover} would imply that
$\pi(G)=\pi(H_{r-1}\oplus_3 Q_r)=4$, contrary to the
assumption. If we use the same argument inductively with
decreasing index, we can conclude that every graph $H_i$ with
$p+1\le i\le r$ has $\pi(H_i)\ge 5$.

Next we prove that each distinguished vertex $u_i$ of $H_i$
with $p\le i\le r-1$ is an apex. Recall that each $H_i$ arises
from $Pg$ by a repeated application of edge-insertion or
vertex-substitution. Since every vertex in the Petersen graph
is an apex, repeated use of Lemma~\ref{lem:pmi4sum3gen} implies
that each $u_i$ for $i\geq p$ is an apex of $H_i$.

Finally we prove that each $Q_i$ with $i\ge p+1$ is
quasi-bipartite. We already know that $\pi(H_i)\ge 5$ and that
$u_i$ is apex for each $i\ge p$. Since $H_{i+1}=H_i\oplus_3
Q_{i+1}$ whenever $p\le i\le r-1$, from
Theorem~\ref{thm:3-sumpmi5-converse} we infer that $Q_{i+1}$ is
quasi-bipartite and $H_{i+1}$ arises from $H_i$ by a correct
substitution of $u_i$. This completes the proof of Claim~2 and
also the proof of the theorem.
\end{proof}

\end{document}
